\par\noindent\begin{minipage}{\linewidth}
\tagstructbegin{tag=Figure,alttext={绝对值任务的三个点为负一一、零零、一一。最佳直线是高度三分之二的水平虚线，无法通过三个点；两个 ReLU 单元相加形成 V 形实线，精确通过三个点。}}
\begin{center}
\begin{tikzpicture}
\begin{axis}[width=112mm,height=47mm,xmin=-1.3,xmax=1.3,ymin=-.1,ymax=1.35,
 axis lines=left,xlabel={$x$},ylabel={预测与目标},xtick={-1,0,1},ytick={0,.6666667,1},
 yticklabels={$0$,$2/3$,$1$},tick label style={font=\footnotesize},
 legend style={at={(.5,1.08)},anchor=south,font=\footnotesize,draw=none,legend columns=2}]
\addplot[flowcolor,thick,dashed,domain=-1.15:1.15]{2/3};
\addlegendentry{最佳直线}
\addplot[flowcolor,very thick] coordinates {(-1.15,1.15)(0,0)(1.15,1.15)};
\addlegendentry{两个 ReLU 相加}
\addplot[only marks,datacolor,mark=*,mark size=2.5pt] coordinates {(-1,1)(0,0)(1,1)};
\end{axis}
\end{tikzpicture}
\captionof{figure}{同一组三点：直线保留误差，两个 ReLU 单元精确拟合}
\label{fig:relu-fit}
\end{center}
\tagstructend
\end{minipage}\par
